\section{工作经历}
\cventry
{研究助理} % Job Title
{深圳河套学院 (SLAI)} % Organization
{中国深圳} % Location
{2026年6月 - 至今} % Date
\begin{cvitems}
\item {从事机器人与具身智能方向的研究工作。}
\end{cvitems}

\section{项目经历}
% \cventry
% {Stable Diffusion图像生成模型微调} % Job Title
% {本科生研究机会(UROP)} % Organization
% {HKUST} % Location
% {2024.2 - 2024.5} % Date
% \begin{cvitems}
% \item {使用LoRA微调Stable Diffusion图像生成模型，实现了改进的风格适应。通过集成ControlNet进行精确条件控制，进一步提升了生成稳定性和可控性。}
% \end{cvitems}

\cventry
{在研项目，由香港科技大学与深圳河套学院联合开展} % Job Title
{乐高装配理解} % Organization
{指导老师：汪子琦教授} % Location
{2026年3月 - 至今} % Date
\begin{cvitems}
\item {研究乐高装配的自动化理解，结合计算机视觉与规划方法，从可视化装配说明与三维模型中还原逐步装配结构。}
\end{cvitems}

\cventry
{课程项目} % Job Title
{ELEC5660 空中机器人导论} % Organization
{指导老师：沈劭劼教授} % Location
{2026年2月 - 2026年5月} % Date
\begin{cvitems}
\item {从零开始实现空中机器人算法，包括视觉惯性里程计（VIO）、扩展卡尔曼滤波（EKF）状态估计以及自主路径规划，在仿真环境中实现了稳健的无人机定位、建图与导航。}
\end{cvitems}

\cventry
{科研实习} % Job Title
{四足机器人控制} % Organization
{指导老师：Maurice Pagnucco教授和Yang Song教授} % Location
{2025年2月 - 2025年8月} % Date
\begin{cvitems}
\item {构建并编程宇树四足机器人以执行复杂任务，集成深度相机进行准确的物体定位（定位误差小于10cm）和激光雷达进行平面建图（SLAM）。}
\end{cvitems}

\cventry
{毕业设计} % Job Title
{水下多相机定位系统} % Organization
{指导老师：尹欢教授和张福民教授} % Location
{2025年6月 - 2026年5月} % Date
\begin{cvitems}
\item {开发了基于多相机视觉的定位系统，用于跟踪多个自主水下航行器（AUV）。降低在水下复杂光线场景下的定位误差，在1m$\times$1m范围内提升定位精度至10cm以内。}
\end{cvitems}


